\documentclass[11pt,a4paper]{article}
%-------------------------
% Basic packages
%-------------------------
\usepackage[utf8]{inputenc}
\usepackage[T1]{fontenc}
\usepackage{lmodern}
\usepackage{microtype}
\usepackage{geometry}
\geometry{margin=1in}

% Mathematics
\usepackage{amsmath,amssymb,amsthm,mathtools,bm}
\usepackage{physics}    % convenient bra-ket, derivatives, etc.
\usepackage{siunitx}    % units
\usepackage{braket}     % alternate bra/ket macros
\usepackage{mhchem}     % chemistry notation if needed

% Figures & graphics
\usepackage{graphicx}
\usepackage{tikz}
\usepackage{pgfplots}
\pgfplotsset{compat=1.18}

% References, hyperlinks, and clever names
\usepackage[colorlinks=true,linkcolor=blue,citecolor=blue,urlcolor=blue]{hyperref}
\usepackage[nameinlink]{cleveref}

% Lists and layout
\usepackage{enumitem}
\setlist{itemsep=2pt,parsep=2pt}

% Theorem environments
\theoremstyle{plain}
\newtheorem{theorem}{Theorem}[section]
\newtheorem{lemma}[theorem]{Lemma}
\newtheorem{proposition}[theorem]{Proposition}
\newtheorem{corollary}[theorem]{Corollary}

\theoremstyle{definition}
\newtheorem{definition}[theorem]{Definition}
\newtheorem{example}[theorem]{Example}

\theoremstyle{remark}
\newtheorem{remark}[theorem]{Remark}

%-------------------------
% Useful macros
%-------------------------
% \newcommand{\dd}{\mathrm{d}}                      % differential
\newcommand{\ii}{\mathrm{i}}                      % imaginary unit
\newcommand{\eexp}{\mathrm{e}}                    % exponential base
\DeclareMathOperator{\sgn}{sgn}

% Shortcuts for common physics notation (overrides safe if physics package present)
\renewcommand{\vec}[1]{\boldsymbol{#1}}

%-------------------------
% Bibliography (biber/biblatex) - optional
%-------------------------
% \usepackage[backend=biber,style=phys]{biblatex}
% \addbibresource{references.bib}

%-------------------------
% Document metadata
%-------------------------
\title{Execise 2}
\author{Zhuge X.K.}
\date{\today}

%-------------------------
% Document
%-------------------------
\begin{document}

\maketitle
\section*{Exercise 1.2(E.2.1, page 112)}
In Section 2.1 we calculate the contact resistance when a narrow conductor with M modes is connect to two very wide contacts. If the number of modes in the contacts is not infinite, but some finite number, N, then the left-moving and right-moving carriers inside the contacts have different electrochemical potentials as shown in Fig. E.2.1 (page 112 of Datta's book).Show that the contact resistance taking this into account is given by
\[
    R_c = \frac{h}{2e^2}\left[\frac{1}{M}-\frac{1}{N}\right]
\]
Assume reflectionless contacts as in the text. For futher discussions on the nature of the contact resistance at different types of interfaces see R.Landauer(1989), hys. Cond. Matter, 1, 8099 and M. C. Payne(1989), J. Phys. Cond. Matter, 1, 4931.

\begin{proof}
Because of the conservation of current, the current flowing through the conductor is given by
\begin{equation}
    I = \frac{2e}{h}M[\mu_L - \mu_R] = \frac{2e}{h}N[\mu_L - \mu'] = \frac{2e}{h}N[\mu' - \mu_R]
\end{equation}
To simplify the algebra set $\mu_L = 1, \mu_R = 0$. This yields
\begin{equation}
    1 - \mu' = \mu'' = \frac{M}{N}
\end{equation}
Hence 
\begin{equation}
    eV = \frac{\mu_L + \mu'}{2} - \frac{\mu'' + \mu_R}{2} = \frac{1 + 1 - \frac{M}{N}}{2} - \frac{\frac{M}{N} + 0}{2} = 1 - \frac{M}{N}
\end{equation}
\begin{equation}
    I = \frac{2e}{h}M
\end{equation}
Thus the contact resistance is given by
\begin{equation}
    R_c = \frac{V}{I} = \frac{h}{2e^2}\left[\frac{1}{M} - \frac{1}{N}\right]
\end{equation}
\end{proof}

\end{document}